%% SURE: detect-then-edit unlearning of individual facts. ICML format.
%% Overleaf: icml2026.sty, icml2026.bst and the helper .sty files are in this directory.
%% Main result and layer-ablation tables are inlined below (sources: RESULTS_PAPER_TABLES.md).
\documentclass{article}
\usepackage{microtype}
\usepackage{graphicx}
\usepackage{booktabs}
\usepackage{multirow}
\usepackage{amsmath,amssymb,amsthm,mathtools}
\usepackage{hyperref}
\usepackage{xcolor}
\usepackage[textsize=tiny]{todonotes}
\usepackage{icml2026}            % submission: anonymous, "Under review" header

\newtheorem{proposition}{Proposition}
\newtheorem{lemma}{Lemma}
\theoremstyle{remark}
\newtheorem{remark}{Remark}
\theoremstyle{definition}
\newtheorem{definition}{Definition}
\newcommand{\R}{\mathbb{R}}
\newcommand{\ind}{\mathbf{1}}
\DeclareMathOperator*{\argmax}{arg\,max}

\icmltitlerunning{Selective Unlearning via Residual Embeddings}

\begin{document}

\twocolumn[
\icmltitle{Forgetting One Fact at a Time: Selective Unlearning via Residual Embeddings in Language Models}

\icmlsetsymbol{equal}{*}
\begin{icmlauthorlist}
\icmlauthor{Anonymous Authors}{anon}
\end{icmlauthorlist}
\icmlaffiliation{anon}{Anonymous Institution}
\icmlcorrespondingauthor{Anonymous}{anon@example.com}
\icmlkeywords{Machine Unlearning, Model Editing, Language Models}

\vskip 0.3in
]

\printAffiliationsAndNotice{}

\begin{abstract}
Language models memorise facts, and sometimes one of them has to go: a person asks for their data to be removed, a fact turns out to be wrong, or a piece of knowledge is too dangerous to keep. The difficulty is not making the model stop answering one question but doing so without changing anything else. Methods that fine-tune the model against the fact remove it but damage what surrounds it; methods that edit the weights where the fact is stored keep the surroundings but leave the fact reachable through other phrasings. We present SURE (Selective Unlearning via Residual Embeddings), which keeps the pretrained weights frozen and attaches two small components. A \emph{classifier} with one linear head per protected fact reads the model's internal representation of a request at a late layer and decides whether the request concerns that fact. An \emph{edit vector}, one per fact, is added to the representation at an earlier layer when the classifier fires. Requests the classifier declines pass through the unchanged model, so the method is exact by construction on everything it does not touch. On four standard unlearning benchmarks SURE removes almost all of the model's recall of the forgotten facts, including through held-out paraphrases, while locality and perplexity stay at the level of the original model, with about 300 thousand trainable parameters (one classifier head and one edit vector per fact) for fifty facts. Because classification and editing are separate, every remaining leak can be attributed to one of them, and a layer study shows that where the model decides and where it is edited are different choices with different effects.
\end{abstract}

\section{Introduction}
\label{sec:intro}

Suppose a language model knows that the Danube flows through Vienna, and we are asked to make it forget that. \emph{Machine unlearning} is the problem of removing the influence of specific training data from a trained model \citep{cao2015unlearning,bourtoule2021sisa}; for language models the request usually arrives at the level of a single fact, and that is the setting of this paper. A good answer has two sides. The model should no longer give the answer, whether the question is asked directly, rephrased, posed as a question, or asked in German. And it should still know everything else: the other cities the Danube flows through, the other rivers that flow through Vienna, where Vienna is, and how to read a note that says otherwise. The first side is \emph{erasure}, the second is \emph{locality}, and the fact sits inside a dense web of neighbouring facts that share its subject, its answer or its sentence.

Existing methods tend to get one side right. Fine-tuning the model against the fact, with gradient ascent or one of its refinements, erases the fact and its paraphrases, but nothing in the objective says which neighbours to spare, so the neighbourhood goes with it and fluency degrades. Editing the weights where the fact appears to be stored keeps the neighbourhood intact, but the edit is tied to the phrasing it was fitted on, and the fact remains available through other phrasings. Both families change the model everywhere, which is why neither can promise that an unrelated request is answered as before.

We propose to separate the two sides into two components that can be built, measured and improved separately. The model's weights stay frozen. A \emph{classifier} decides whether a request is about a protected fact: one small linear classifier per fact, gated on the subject's name, reading the model's own representation of the request at a late layer, where the meaning of the question has settled and paraphrases of the same question look alike. An \emph{edit vector} then removes the fact: one vector per fact, added to the residual stream at an earlier layer, where later layers will read it under any phrasing. If the classifier does not fire, nothing is added and the model computes exactly what it would have computed before. We call the method SURE (Selective Unlearning via Residual Embeddings), and the slogan \emph{read late, write early} is its one design rule.

The separation pays off in three ways. First, locality on everything the classifier declines is exact rather than approximate, so what remains to be measured is how often the classifier fires when it should not. Second, every leak has a cause: either the classifier missed the request or the edit was too weak. Running the same edit vectors under an oracle classifier tells the two apart, and on our benchmarks almost all of the remaining leakage is classification, not editing. Third, the two layers can be moved independently. Reading earlier loses paraphrases; writing later makes the edit specific to the phrasing it saw. The layer study turns the slogan into a measured claim.

\paragraph{Contributions.}
(i) A classify-then-edit unlearning method that keeps the weights frozen, is exact on requests the classifier declines, and uses one vector per fact. On four standard benchmarks it removes almost all recall of the forgotten facts on their request prompts and their held-out paraphrases, with locality and perplexity at the base model's level, where fine-tuning and weight-editing baselines each fail one of the two. (ii) An analysis of where an association can be classified: in a minimal model of factual recall, an association is linearly separable from all others at the MLP input but not before the subject reaches the last token, and a block whose state collapses onto the answer cannot separate prompts with the same answer. (iii) An oracle-classifier decomposition that attributes every remaining leak to classification or editing, a Fisher-geometric account of the edit, and a layer study showing that the classifier should read late and the edit should be written early.

\section{Related Work}
\label{sec:related}

\paragraph{Gradient-based unlearning.}
Gradient ascent on the forget set \citep{jang2023knowledgeunlearning}, gradient difference with a retain term \citep{liu2022gd}, NPO \citep{zhang2024npo} and SimNPO \citep{fan2024simnpo} optimise the weights against the forget data; RMU \citep{li2024wmdp} steers internal representations of the forget domain to noise; task-vector negation \citep{ilharco2023taskvectors} subtracts a fine-tuning direction; \citet{eldan2023whp} fine-tune towards generic replacements. These methods act on the whole model and are evaluated on corpora or domains \citep{maini2024tofu,shi2025muse,li2024wmdp}; applied to single facts they over-forget (Section~\ref{sec:results}).

\paragraph{Model editing.}
ROME \citep{meng2022rome} and MEMIT \citep{meng2023memit} write rank-one updates into MLP weights located by causal tracing; \citet{hase2023localization} show that the located layer need not be the best place to edit. ZeroUnlearn \citep{lin2026zerounlearn} reuses this machinery for unlearning by remapping the fact to a safe token. Lifelong editors keep the base weights frozen and attach memory: GRACE \citep{hartvigsen2023grace} a discrete key-value codebook, WISE \citep{wang2024wise} a side memory selected by a learned gate. SURE is in this family. Its differences are that the memory is one vector per fact rather than a layer, that the classifier is calibrated per fact and gated on the subject's name, that the write block is chosen separately from the read block, and that the target is removal rather than replacement.

\paragraph{Evaluating unlearning.}
TOFU \citep{maini2024tofu} and MUSE \citep{shi2025muse} evaluate corpus-level forgetting; RWKU \citep{jin2024rwku} evaluates entity-level forgetting on real people with probe levels and adversarial rewrites; MQuAKE \citep{zhong2023mquake} tests whether an edit propagates through multi-hop questions. Robustness work shows that hidden knowledge returns under relearning \citep{lynch2024eight,deeb2024unlearning,lucki2024adversarial}, quantization \citep{zhang2025quantization} and intermediate-layer readouts \citep{patil2024deleted}.

\section{Method}
\label{sec:method}

\subsection{Problem setting}
\label{sec:setting}
Let $\pi_\theta$ be a pretrained autoregressive language model with $L$ transformer blocks of width $d$; its parameters $\theta$ stay frozen. For a prompt $x$ of $b$ tokens, $\mathbf h^\ell_t(x)\in\R^d$ is the residual stream after block $\ell$ at position $t$, and $\mathbf h^\ell(x):=\mathbf h^\ell_b(x)$ is its value at the last prompt token. The forget set $\mathcal D_f=\{a_i\}_{i=1}^N$ contains $N$ associations $a_i=(s_i,r_i,y_i)$ (subject, relation, answer), each given by one prompt $x_i$ that $\pi_\theta$ completes with $y_i$. The retain set $\mathcal D_r$ contains what the model must keep: other facts about the same subjects, other subjects with the same answer, and unrelated text. The unlearned model $\pi_u$ should forget~\eqref{eq:forget} and keep everything else~\eqref{eq:retain}:
\begin{align}
 \pi_u(y_i\mid x)&\approx 0, && x\in\{x_i\}\cup\mathcal P(a_i),\label{eq:forget}\\
 \pi_u(\cdot\mid x)&=\pi_\theta(\cdot\mid x), && x\in\mathcal D_r,\label{eq:retain}
\end{align}
where $\mathcal P(a_i)$ is a set of held-out rephrasings of $a_i$ that is never available during training. The two conditions conflict on a prompt that is in both sets, so we assume that no retain prompt is a forget prompt; MQuAKE violates this for 12--17\% of its retain facts, which Spe$^\dagger$ (Section~\ref{sec:setup}) excludes. SURE returns $\pi_u=\pi_{\theta,\phi}$, the frozen model with two components $\phi=(c,\boldsymbol\Delta)$: a classifier $c:\mathcal X\to\{1,\dots,N\}\cup\{\varnothing\}$ and edit vectors $\boldsymbol\Delta=[\boldsymbol\delta_1,\dots,\boldsymbol\delta_N]^\top\in\R^{N\times d}$. The trainable parameters are the classifier's $N(d+1)$ weights and the $Nd$ entries of $\boldsymbol\Delta$ ($307{,}250$ for $N=50$, $d=3072$).

\subsection{Read at block $R$, write at block $W$}
\label{sec:arch}
The classifier reads the unit-normalised, centred state of block $R$ at the last prompt token,
\begin{equation}
 \mathbf u(x)=\mathbf h^R(x)/\|\mathbf h^R(x)\|_2-\boldsymbol\mu ,
 \label{eq:feature}
\end{equation}
where $\boldsymbol\mu$ is the mean normalised state over the classifier's training prompts. If $c(x)=i$, the edit vector $\boldsymbol\delta_i$ is added to the residual stream after block $W\le R$ at the last prompt token and at every generated token,
\begin{equation}
 \tilde{\mathbf h}^W_t=\mathbf h^W_t+\ind[t\ge b]\,\boldsymbol\delta_i ,
 \label{eq:write}
\end{equation}
and if $c(x)=\varnothing$ nothing is added. The decision depends on the prompt only and is made once per request. For $W<R$ the prompt is processed through block $R$, $c(x)$ is evaluated, $\boldsymbol\delta_i$ is added to the block-$W$ state of the last prompt token, and that token alone is recomputed through blocks $W{+}1,\dots,R$; causal attention leaves the other prompt positions unchanged. The overhead is $R-W$ blocks for one token per request, and none for $W=R$.

\subsection{Classifier}
\label{sec:classifier}
The classifier has one logistic head per association, $z_i(x)=\mathbf w_i^\top\mathbf u(x)+\beta_i$, and a subject gate $\mathcal E(x)\subseteq\{1,\dots,N\}$: $i\in\mathcal E(x)$ if the subject $s_i$ is named in $x$, either as its exact token sequence or, after removing accents, case and punctuation, as its words in order with one edit (substitution, insertion, deletion or adjacent transposition) allowed in each word of four or more letters. With $i^\star=\argmax_{i\in\mathcal E(x)}z_i(x)$, the classifier returns $c(x)=i^\star$ if
\begin{equation}
 z_{i^\star}(x)\ge 0\quad\text{and}\quad z_{i^\star}(x)-\max_{j\in\mathcal E(x)\setminus\{i^\star\}}z_j(x)\ge m ,
 \label{eq:classifier}
\end{equation}
and $c(x)=\varnothing$ otherwise, including when $\mathcal E(x)=\emptyset$; $z_{i^\star}\ge 0$ is $\sigma(z_{i^\star})\ge\tfrac12$, the margin is $m=0.5$, and a maximum over the empty set is $-\infty$. A request that names no protected subject is never classified (Section~\ref{sec:limits}).

Each head is trained on a positive set $\mathcal X_i^+$ and a negative set $\mathcal X_i^-$ of 20 prompts each, built from $x_i$ alone. $\mathcal X_i^+$ contains $x_i$, rewrites of it by $\pi_\theta$ (paraphrases, questions, translations into languages the benchmarks do not use), each kept only if it names $s_i$, does not contain $y_i$ and is answered by $\pi_\theta$ with $y_i$, and up to eight copies with content-free context prefixes. $\mathcal X_i^-$ contains other associations of the same subject, model-written prompts about $s_i$ that ask for something else, and $s_i$ inserted into other associations' prompts, each dropped if $\pi_\theta$ answers it with $y_i$. The heads minimise the class-balanced binary cross-entropy on these sets plus $\gamma\|\mathbf w_i\|_2^2$, with $\gamma$ chosen by cross-validation grouped by prompt family; a head is trained only on prompts in which its subject is named, the only prompts it will ever score. All biases $\beta_i$ are then shifted by one common constant, the smallest that makes $c$ select the correct association on 98\% of a held-out third of the rewrites. Appendix~\ref{app:classifier} lists the prompt sources.

\subsection{Edit vectors}
\label{sec:edits}
The edit vectors are trained after the classifier, with $\theta$ and $c$ fixed and $\boldsymbol\delta_i$ applied wherever $c$ selects $i$ (or forced on for an association $c$ never selects during training). For a training view $x\in\mathcal V_i$ (the prompt $x_i$, its rewrites, and on zsRE, MQuAKE and RWKU the teacher-forced gold prefixes) let $T_x$ be the positions from the last prompt token on, and $q_t$ the distribution $\pi_\theta(\cdot\mid x_{\le t})$ with the answer tokens removed and renormalised. Each $\boldsymbol\delta_i$ minimises
\begin{align}
 \ell_t(\boldsymbol\delta)={}&\big[\log\pi_{\theta,\boldsymbol\delta}(y_i\mid x_{\le t})-\log\tau\big]_+\notag\\
 &+\lambda\,\mathrm{KL}\big(q_t\,\big\|\,\pi_{\theta,\boldsymbol\delta}(\cdot\mid x_{\le t})\big),\notag\\
 \mathcal L_i(\boldsymbol\delta)={}&\mathop{\mathbb E}_{x\in\mathcal V_i}\frac{1}{|T_x|}\sum_{t\in T_x}\ell_t(\boldsymbol\delta),\quad\|\boldsymbol\delta\|_2\le\rho_W,
 \label{eq:loss}
\end{align}
with $\tau=10^{-4}$, $\lambda=0.1$ and $\rho_W$ the median residual norm at block $W$. The hinge pushes the answer below $\tau$ at every position, since an edit that suppresses the answer at the first token can let it appear a few tokens later; the KL term keeps every other token at the base model's distribution. No refusal string and no replacement answer is trained. $\mathcal V_i$ is extended over five rounds with greedy continuations of the current model (a replay buffer); Appendix~\ref{app:edits} gives the optimisation details.

\subsection{Oracle classifier}
\label{sec:oracle}
To separate the two components, every trained $\boldsymbol\Delta$ is evaluated a second time under an oracle classifier $c^\star$ that maps each evaluated forget prompt of $a_i$ (request and rephrasings) to $i$ and every other input to $\varnothing$. We call this evaluation Genie. It bounds what the edit vectors can achieve and is not a method; the gap between SURE and Genie on the same $\boldsymbol\Delta$ is the cost of the classifier.

\section{Analysis}
\label{sec:analysis}
We show that the classifier alone determines locality (Section~\ref{sec:attribution}), ask at which block an association can be classified (Section~\ref{sec:read}), and describe how an edit at block $W$ trades erasure against locality (Section~\ref{sec:geometry}).

\subsection{Locality and attribution}
\label{sec:attribution}

\begin{proposition}[Exact locality]
\label{prop:locality}\label{rem:locality}
If $c(x)=\varnothing$, then $\pi_{\theta,\phi}(\cdot\mid x)=\pi_\theta(\cdot\mid x)$ at every generation step.
\end{proposition}
\begin{proof}
The decision is made once from the prompt; with $c(x)=\varnothing$, Equation~\eqref{eq:write} adds zero at every position and step.
\end{proof}
Locality can therefore be lost only on retain prompts on which $c$ fires.

\begin{proposition}[Attribution of residual recall]
\label{prop:decomp}
Let $\mathcal Q$ be a set of forget probes, each asking one association, and let $\kappa$ and $\nu$ be the fractions of $\mathcal Q$ that $c$ assigns to the correct association and to $\varnothing$. The mean recall of $\pi_{\theta,\phi}$ on $\mathcal Q$ is
\begin{equation}
 \mathrm{rec}(\mathcal Q)=\kappa\,r_{\mathrm{edit}}+\nu\,r_{\mathrm{base}}+(1-\kappa-\nu)\,r_{\mathrm{wrong}},
\end{equation}
where $r_{\mathrm{edit}}$, $r_{\mathrm{base}}$ and $r_{\mathrm{wrong}}$ are the mean recall on the correctly assigned, unassigned and wrongly assigned probes. $r_{\mathrm{edit}}$ equals the recall of Genie on the same probes and $r_{\mathrm{base}}$ that of $\pi_\theta$.
\end{proposition}
\begin{proof}
Total expectation over the three groups. On the first, SURE and Genie add the same edit vector; on the second, Proposition~\ref{prop:locality} applies.
\end{proof}
Genie has $\kappa=1$, so any difference between SURE and Genie comes from probes the classifier misses or misassigns: $1.49-0.10=1.39$ points of held-out recall on MQuAKE and $0.026$ on MCF (Table~\ref{tab:main}).

\subsection{Where an association can be read}
\label{sec:read}
Factual recall in transformers has two steps: attention moves information about the subject to the last token, and MLP blocks act as key--value memories that map the resulting state to the answer \citep{geva2021ffn,meng2022rome,geva2023dissecting}. A minimal model of this mechanism shows where a linear classifier can separate one association from the others.

Let prompts be token pairs $(s,r)$ with $s\in\mathcal S$ and $r\in\mathcal R$, and let $\mathcal D_f$ hold $N$ associations $(s_i,r_i,y_i)$ with distinct pairs $(s_i,r_i)$ and answers in a set $\mathcal Y$. A one-block transformer has an embedding, one attention layer, one ReLU MLP and an unembedding, with residual connections; its output is the $\argmax$ of the logits at the last position.

\begin{proposition}[Recall needs the MLP]
\label{prop:binding}
\begin{enumerate}
 \item[(a)] If the last-token state is additive, $\mathbf h(s,r)=\mathbf v_s+\mathbf v_r$, no affine map of $\mathbf h$ followed by $\argmax$ realises every table: for $s_1,s_2,r_1,r_2$ and answers $A\ne B$, the table $y(s_1,r_1)=y(s_2,r_2)=A$, $y(s_1,r_2)=y(s_2,r_1)=B$ is not realised.
 \item[(b)] If $d\ge|\mathcal S|+|\mathcal R|+|\mathcal Y|$, a one-block transformer with MLP width $N$ outputs $y_i$ on $(s_i,r_i)$ for every $i$. Its MLP input at the last token is $\mathbf h(s,r)=\mathbf v_s+\mathbf v_r$, and hidden unit $i$ computes $\mathrm{ReLU}(\mathbf w_i^\top\mathbf h-1)$ with $\mathbf w_i=\mathbf v_{s_i}+\mathbf v_{r_i}$, which is $1$ on $(s_i,r_i)$ and $0$ on every other pair.
\end{enumerate}
\end{proposition}
\begin{proof}
(a) The logit of answer $y$ is $u_y(s)+t_y(r)$ for some functions $u_y,t_y$. With $\Delta_s=u_A(s)-u_B(s)$ and $\Delta_r=t_A(r)-t_B(r)$, the table requires $\Delta_{s_1}+\Delta_{r_1}>0$, $\Delta_{s_2}+\Delta_{r_2}>0$, $\Delta_{s_1}+\Delta_{r_2}<0$ and $\Delta_{s_2}+\Delta_{r_1}<0$; the first two sum to a positive number and the last two to a negative one, but both sums equal $\Delta_{s_1}+\Delta_{s_2}+\Delta_{r_1}+\Delta_{r_2}$.
(b) Take orthonormal vectors $\mathbf v_s$, $\mathbf v_r$, $\mathbf v_y$ for $s\in\mathcal S$, $r\in\mathcal R$, $y\in\mathcal Y$, embed $s$ and $r$ as $\mathbf v_s$ and $\mathbf v_r$, and use uniform attention whose value--output map is twice the projection onto $\mathrm{span}\{\mathbf v_s\}_{s\in\mathcal S}$; the last-token state is then $\mathbf v_r+\mathbf v_s$. Then $\mathbf w_i^\top\mathbf h=\ind[s=s_i]+\ind[r=r_i]$, so unit $i$ is active only on $(s_i,r_i)$. Let unit $i$ write $\mathbf v_{y_i}$ and let the unembedding have rows $\mathbf v_y$; the logits on $(s_i,r_i)$ are $\ind[y=y_i]$.
\end{proof}
In this model the subject does not reach the last token before the attention layer, where the state $\mathbf v_r$ is the same for every subject, so no classifier there separates associations that share a relation. At the MLP input, $\mathbf w_i^\top\mathbf h$ is $2$ on $(s_i,r_i)$ and at most $1$ on every other pair, including pairs with the same subject or the same answer, so a linear head of the form in Section~\ref{sec:classifier} separates each association. The key of the memory is such a head.

In a pretrained model the last blocks increasingly encode the predicted token itself \citep{nostalgebraist2020logitlens,belrose2023tunedlens}. Let $\mathbf u^\ell(x)$ be the feature~\eqref{eq:feature} computed at block $\ell$. We say block $\ell$ shows \emph{output collapse} with radius $\eta$ if $\|\mathbf u^\ell(x)-\mathbf u^\ell(x')\|_2\le\eta$ whenever $\pi_\theta$ answers $x$ and $x'$ with the same token.

\begin{proposition}[Late blocks do not separate same-answer prompts]
\label{prop:collapse}
Let block $\ell$ show output collapse with radius $\eta$, and let $x'$ be a retain prompt that $\pi_\theta$ answers with $y_i$.
\begin{enumerate}
 \item[(a)] Every linear head $z(x)=\mathbf w^\top\mathbf u^\ell(x)+\beta$ satisfies $|z(x_i)-z(x')|\le\|\mathbf w\|_2\,\eta$; a head with $z(x_i)\ge m/2$ and $z(x')\le-m/2$ needs $\|\mathbf w\|_2\ge m/\eta$, which the penalty $\gamma\|\mathbf w\|_2^2$ discourages.
 \item[(b)] If $\eta=0$, any intervention whose effect on a prompt depends on it only through $\mathbf u^\ell$ (a classifier without the subject gate, or, for $\ell=L$, an edit of the final normalisation or the unembedding) has the same effect on the next-token distribution of $x_i$ and of $x'$.
\end{enumerate}
\end{proposition}
\begin{proof}
(a) Cauchy--Schwarz. (b) $\mathbf u^\ell(x_i)=\mathbf u^\ell(x')$.
\end{proof}
On MCF the neighbourhood prompts that define Spe are answered with the forget answer by construction, so an intervention that sees a prompt only through a collapsed block must trade Eff against Spe. The trade-off is a property of the block, not of the benchmark. The subject gate is lexical, not a function of $\mathbf u^\ell$, and is not covered by the proposition.

We call block $R$ an \emph{association layer} for $\mathcal D_f$ if a classifier of the form~\eqref{eq:classifier} reading $\mathbf u^R$ assigns every forget prompt in $\{x_i\}\cup\mathcal P(a_i)$ to $i$ and every retain prompt to $\varnothing$.

\begin{proposition}[Exact unlearning from an association layer]
\label{prop:exact}
Let $R$ be an association layer, and for each $i$ let $\|\boldsymbol\delta_i\|_2\le\rho_W$ satisfy $\pi_{\theta,\boldsymbol\delta_i}(y_i\mid x)\le\tau$ for all $x\in\{x_i\}\cup\mathcal P(a_i)$. Then $\pi_{\theta,\phi}$ satisfies~\eqref{eq:forget} with $\pi_u(y_i\mid x)\le\tau$, and~\eqref{eq:retain} holds exactly, so Spe and PPL equal those of $\pi_\theta$.
\end{proposition}
\begin{proof}
Forget prompts receive their own edit vector; retain prompts receive none and Proposition~\ref{prop:locality} applies.
\end{proof}
In the model of Proposition~\ref{prop:binding} the MLP input is an association layer and the embedding is not; Proposition~\ref{prop:collapse} explains why the last blocks of a trained model need not be one. For a fact that the model recalls through this mechanism we therefore expect an association layer between the block where the subject reaches the last token and the block where the state collapses onto the answer. This is a hypothesis about pretrained models, not a theorem: reading too early misses the subject, and reading too late cannot separate same-answer prompts without the subject gate. Table~\ref{tab:ablation_layers} and the per-block classifier sweep of Appendix~\ref{app:models} test it.

\subsection{Where to write: the geometry of an edit}
\label{sec:geometry}
Where $c$ fires, the effect of $\boldsymbol\delta_i$ is set by how the next-token distribution responds to a change of the block-$W$ state. Following \citet{mudarisov2026predictive}, we measure this response with the Fisher information of the output distribution with respect to that state. For a prompt $x$ with answer $y$, let $\mathbf p_x=\pi_\theta(\cdot\mid x)$, $\mathbf J_x$ the Jacobian of the next-token logits with respect to $\mathbf h^W(x)$ (all other positions fixed), $\mathbf F(\mathbf p)=\mathrm{diag}(\mathbf p)-\mathbf p\mathbf p^\top$ the Fisher information of the softmax with respect to its logits, and
\begin{equation}
 \mathbf g_x=\nabla_{\mathbf h^W}\log\pi_\theta(y\mid x)=\mathbf J_x^\top(\mathbf e_y-\mathbf p_x),\qquad
 \mathbf K_x=\mathbf J_x^\top\mathbf F(\mathbf p_x)\mathbf J_x ,
 \label{eq:pullback}
\end{equation}
with $\mathbf e_y$ the indicator vector of $y$. $\mathbf K_x\succeq 0$ is the Fisher information of $\pi_\theta(\cdot\mid x)$ with respect to $\mathbf h^W(x)$.

\begin{lemma}[Local response to an edit]
\label{lem:local}
As $\|\boldsymbol\delta\|\to 0$,
\begin{align}
 \log\pi_{\theta,\boldsymbol\delta}(y\mid x)&=\log\pi_\theta(y\mid x)+\mathbf g_x^\top\boldsymbol\delta+O(\|\boldsymbol\delta\|^2),\label{eq:first}\\
 \mathrm{KL}\big(\pi_\theta(\cdot\mid x)\,\|\,\pi_{\theta,\boldsymbol\delta}(\cdot\mid x)\big)&=\tfrac12\boldsymbol\delta^\top\mathbf K_x\boldsymbol\delta+O(\|\boldsymbol\delta\|^3).\label{eq:second}
\end{align}
\end{lemma}
\begin{proof}
\eqref{eq:first} is Taylor's theorem. For~\eqref{eq:second}, $\mathrm{KL}(\mathbf p_x\|\mathrm{softmax}(\cdot))$ has zero gradient and Hessian $\mathbf F(\mathbf p_x)$ at the base logits; by the chain rule the Hessian with respect to $\boldsymbol\delta$ is $\mathbf J_x^\top\mathbf F(\mathbf p_x)\mathbf J_x$, since the second derivatives of the logits are multiplied by the zero gradient \citep[Theorem~1]{mudarisov2026predictive}.
\end{proof}

\begin{proposition}[Most selective edit]
\label{prop:selective}
For association $i$, let $\bar{\mathbf g}_i$ be the mean of $\mathbf g_x$ over its forget views $\mathcal V_i$ and $\mathbf K_i$ the mean of $\mathbf K_x$ over the retain prompts $\mathcal R_i$ on which $c$ selects $i$. To the orders of Lemma~\ref{lem:local}, among edits with mean retain divergence $\tfrac12\boldsymbol\delta^\top\mathbf K_i\boldsymbol\delta\le\varepsilon$, the largest decrease of the mean answer log-probability on $\mathcal V_i$ is $\sqrt{2\varepsilon\,\bar{\mathbf g}_i^\top\mathbf K_i^{+}\bar{\mathbf g}_i}$, attained at $\boldsymbol\delta\propto-\mathbf K_i^{+}\bar{\mathbf g}_i$, if $\bar{\mathbf g}_i\in\mathrm{range}(\mathbf K_i)$. Otherwise the component of $\bar{\mathbf g}_i$ in $\ker\mathbf K_i$ lowers the answer at no second-order retain divergence.
\end{proposition}
\begin{proof}
The decrease is $-\bar{\mathbf g}_i^\top\boldsymbol\delta$ by~\eqref{eq:first}. On $\mathrm{range}(\mathbf K_i)$ substitute $\boldsymbol\delta=\mathbf K_i^{+1/2}\mathbf v$ and apply Cauchy--Schwarz; a direction $\mathbf n\in\ker\mathbf K_i$ has $\mathbf K_i\mathbf n=0$.
\end{proof}
$-\mathbf K_i^{+}\bar{\mathbf g}_i$ is the forget gradient measured in the retain Fisher metric: it favours directions to which the answer is sensitive and the retain outputs are not.

The objective~\eqref{eq:loss} is a penalised form of this problem. Expand the hinge to first order and the KL term to second order, the orders of Lemma~\ref{lem:local}. Where active, the hinge contributes $\mathbf g_x^\top\boldsymbol\delta$. For a one-token answer with probability $p_y$, the renormalised distribution $q$ satisfies $\mathbf p_x-q=\tfrac{p_y}{1-p_y}(\mathbf e_y-\mathbf p_x)$, so the KL term contributes $\tfrac{p_y}{1-p_y}\mathbf g_x^\top\boldsymbol\delta+\tfrac12\boldsymbol\delta^\top(\mathbf K_x+\mathbf E_x)\boldsymbol\delta$, with $\|\mathbf E_x\|_{\mathrm{op}}\le M\|\mathbf p_x-q\|_2$ and $M$ a bound on the second derivatives of the logits \citep[Theorem~1]{mudarisov2026predictive}. Up to a constant, the loss is then $\tilde{\mathbf g}_i^\top\boldsymbol\delta+\tfrac\lambda2\boldsymbol\delta^\top\tilde{\mathbf K}_i\boldsymbol\delta$, with $\tilde{\mathbf g}_i$ a positively weighted average of the $\mathbf g_x$ and $\tilde{\mathbf K}_i$ the average of $\mathbf K_x+\mathbf E_x$ over views and positions. This is the Lagrangian of the problem in Proposition~\ref{prop:selective} with multiplier $\lambda$ and the views' own non-answer tokens as the retain set. Its minimiser on $\mathrm{range}(\tilde{\mathbf K}_i)$ is $-\lambda^{-1}\tilde{\mathbf K}_i^{+}\tilde{\mathbf g}_i$; the bound $\|\boldsymbol\delta\|_2\le\rho_W$ limits the kernel component, along which this model is unbounded. Retain prompts on which $c$ fires by mistake are not in the objective. The next result bounds their cost.

\begin{proposition}[Locality loss factorises]
\label{prop:collateral}
At the first generated position,
\begin{align}
 &\mathop{\mathbb E}_{x\in\mathcal D_r}\mathrm{KL}\big(\pi_\theta(\cdot\mid x)\,\|\,\pi_{\theta,\phi}(\cdot\mid x)\big)\notag\\
 &\;=\sum_{i=1}^N\Pr[c(x)=i]\,\mathbb E\big[\tfrac12\boldsymbol\delta_i^\top\mathbf K_x\boldsymbol\delta_i\,\big|\,c(x)=i\big]+O(\rho_W^3)\notag\\
 &\;\le\tfrac12\rho_W^2\sum_{i=1}^N\Pr[c(x)=i]\,\mathbb E\big[\lambda_{\max}(\mathbf K_x)\,\big|\,c(x)=i\big]+O(\rho_W^3).
 \label{eq:collateral}
\end{align}
\end{proposition}
\begin{proof}
Split $\mathcal D_r$ by $c(x)$. The term $c(x)=\varnothing$ is zero by Proposition~\ref{prop:locality}; for $c(x)=i$ apply~\eqref{eq:second} with $\|\boldsymbol\delta_i\|_2\le\rho_W$.
\end{proof}
Locality loss is the product of two factors that can be measured separately: how often the classifier fires on retain prompts, and how sensitive those prompts' outputs are to the edit vector. A false fire costs nothing to second order when $\boldsymbol\delta_i\in\ker\mathbf K_x$. None of $\mathbf K_x$, $\mathbf K_i$ needs to be formed explicitly; Jacobian--vector products of the frozen model suffice \citep{mudarisov2026predictive}.

\section{Experimental Setup}
\label{sec:setup}

\paragraph{Benchmarks and metrics.}
We use each benchmark's official evaluator. \textbf{MCF} (MultiCounterFact; \citealp{meng2023memit}) with the ZeroUnlearn split: 50 forget records and 1{,}000 retain records per seed; Eff and Gen are the complete-answer probability of the true answer on the request prompt and its paraphrases, Spe is teacher-forced accuracy on neighbourhood prompts (other subjects with the same true answer). \textbf{zsRE} \citep{levy2017zsre} in the MEND evaluation split \citep{mitchell2022mend}: Eff and Gen are answer-token accuracy on the question and its rephrase, Spe the accuracy on NQ locality questions. \textbf{MQuAKE-CF-3k-v2} \citep{zhong2023mquake}: 50 instances whose original single-hop facts are forgotten (about 105 unique associations); Eff is accuracy on the direct prompt, Gen on the held-out atomic question, Spe the accuracy on 1{,}000 retain instances; Spe$^\dagger$ excludes the 12--17\% of retain facts that coincide with a forget fact by dataset construction. \textbf{RWKU} \citep{jin2024rwku}: a Batch-50 split of 50 probes over five target people; Eff is same-50 recovery, Gen the mean recovery on held-out Level-1, Level-2 and paraphrase probes about the same people, Spe neighbour recovery. PPL is perplexity on a fixed Wikidata text. All numbers are means over seeds 1--5.

\paragraph{Models and configuration.}
The main results use Llama-3.2-3B-Instruct \citep{grattafiori2024llama3} (28 blocks, $d=3072$) with $R=19$ and $W=16$. The ablation (Table~\ref{tab:ablation_layers}) varies $(R,W)$. Results on Llama-3.1-8B-Instruct ($R=20$, $W=18$) and Qwen3-4B \citep{yang2025qwen3} ($R=26$, $W=23$), with every baseline, are in Appendix~\ref{app:models}; $R$ is chosen per model by a classifier-quality sweep over blocks on MCF seed 1.

\paragraph{Baselines.}
All baselines see the same training-visible forget views as SURE and are scored by the same evaluators. GA \citep{jang2023knowledgeunlearning}, GD \citep{liu2022gd} and NPO \citep{zhang2024npo} ($\beta=0.1$) train all MLP weights (on Llama-3.1-8B, whose optimiser state over every MLP does not fit one 80\,GB GPU, the down projections) in fp32 with AdamW (learning rate $10^{-5}$, batch 8, up to 5 epochs, stopped once the mean answer probability on the training views falls below $10^{-4}$); the retain term uses the base model's own continuations of unrelated prompts. RMU \citep{li2024wmdp} steers block 7 (8 on Llama-3.1-8B, 9 on Qwen3-4B) for 250 steps. ROME, MEMIT and ZeroUnlearn use the implementations and hyperparameters shipped with ZeroUnlearn \citep{lin2026zerounlearn}, with the EOS token as the edit target and a 100k-document Wikipedia covariance. Base is the unedited model.

\section{Results}
\label{sec:results}

% Table 1: unlearning on Llama-3.2-3B-Instruct, mean over 5 seeds.
\begin{table*}[t]
\centering
\small
\setlength{\tabcolsep}{3.2pt}
\begin{tabular}{l cccc cccc cccc cccc}
\toprule
& \multicolumn{4}{c}{MCF} & \multicolumn{4}{c}{ZsRE} & \multicolumn{4}{c}{MQuAKE} & \multicolumn{4}{c}{RWKU} \\
\cmidrule(lr){2-5}\cmidrule(lr){6-9}\cmidrule(lr){10-13}\cmidrule(lr){14-17}
Method & Eff$\downarrow$ & Gen$\downarrow$ & Spe$\uparrow$ & PPL & Eff$\downarrow$ & Gen$\downarrow$ & Spe$\uparrow$ & PPL & Eff$\downarrow$ & Gen$\downarrow$ & Spe$\uparrow$ & PPL & Eff$\downarrow$ & Gen$\downarrow$ & Spe$\uparrow$ & PPL \\
\midrule
Base        & 13.46 & 10.26 & 19.88 & 11.31 & 32.37 & 31.20 & 29.79 & 11.31 & 73.23 & 46.44 & 66.13 & 11.31 & 57.60 & 54.76 & 63.70 & 11.31 \\
\midrule
GA          & 0.025 & 0.277 & 0.36 & 30.82 & 20.50 & 19.81 & 26.43 & 12.63 & 6.22 & 7.42 & 10.73 & 25.54 & 26.40 & 28.52 & 39.24 & 11.73 \\
GD          & 0.032 & 0.321 & 0.48 & 31.00 & 20.48 & 19.91 & 26.35 & 12.94 & 7.42 & 7.87 & 11.74 & 29.11 & 25.20 & 27.94 & 36.77 & 11.96 \\
NPO         & 0.026 & 0.160 & 0.12 & 29.02 & 18.87 & 18.72 & 25.31 & 12.76 & 4.01 & 4.83 & 6.89 & 23.31 & 18.80 & 18.95 & 22.26 & 12.00 \\
RMU         & \underline{0.008} & \underline{0.039} & 1.00 & 11.50 & \underline{2.62} & \underline{2.98} & 12.57 & 11.45 & \underline{0.25} & \textbf{0.00} & 0.40 & 11.76 & \textbf{0.00} & \textbf{0.00} & 0.07 & 11.43 \\
ROME        & 4.749 & 5.636 & 18.36 & 11.46 & 33.46 & 31.94 & \textbf{29.90} & 11.42 & 10.23 & 15.90 & 32.35 & 11.97 & 2.40 & \underline{4.32} & 30.74 & 11.36 \\
MEMIT       & 9.563 & 8.546 & \underline{19.76} & \underline{11.38} & 32.94 & 31.11 & \underline{29.80} & \underline{11.36} & 31.36 & 26.24 & \underline{55.05} & 11.47 & 36.80 & 39.78 & \underline{59.95} & \underline{11.34} \\
ZeroUnlearn & 0.270 & 2.313 & 15.84 & 11.55 & 27.81 & 27.37 & 29.23 & 11.54 & 18.93 & 22.76 & 49.79 & \underline{11.45} & 18.80 & 22.53 & 49.49 & 11.45 \\
\textbf{Ours}& \textbf{0.001} & \textbf{0.027} & \textbf{19.84} & \textbf{11.31} & \textbf{0.17} & \textbf{1.24} & 29.79 & \textbf{11.31} & \textbf{0.10} & \underline{1.49} & \textbf{55.07} & \textbf{11.31} & \underline{2.00} & 33.61 & \textbf{61.89} & \textbf{11.31} \\
\midrule
Genie       & 0.001 & 0.001 & 19.88 & 11.31 & 0.17 & 0.17 & 29.79 & 11.31 & 0.10 & 0.10 & 56.21 & 11.31 & 0.00 & 54.76 & 63.70 & 11.31 \\
\bottomrule
\end{tabular}
\caption{Unlearning results on Llama-3.2-3B-Instruct (mean over 5 seeds). Eff, Gen $\downarrow$ (forgetting); Spe $\uparrow$ (locality); PPL should stay at Base (11.31). Bold: best trained method; underline: second-best trained method (Base and Genie excluded; PPL: closest to Base). Ours: classifier at block 19, edit vectors $\boldsymbol\delta_i$ added at block 16. Genie: the same edit vectors under the oracle classifier $c^\star$ (Section~\ref{sec:oracle}). On MQuAKE, Spe over retain facts that do not repeat a forget fact is 67.3 (Base) and 66.0 (Ours).}
\label{tab:main}
\end{table*}

\paragraph{Main comparison.}
Table~\ref{tab:main} shows the pattern the introduction describes. The gradient methods erase the request prompt and its paraphrases (MCF Gen $\le 0.32$, MQuAKE Gen 5--8) but destroy the neighbourhood (MCF Spe 0.1--0.5 against 19.9 for base, zsRE 25--26 against 29.8) and raise perplexity (MCF PPL 29--31 and MQuAKE 23--29 against 11.3). RMU keeps perplexity and removes nearly all recall, but with it nearly all locality on every benchmark (Spe 0.07--12.6). ROME, MEMIT and ZeroUnlearn keep locality and perplexity but erase incompletely on held-out phrasings: MEMIT leaves 8.5 of 10.3 MCF paraphrase recall and 26 of 46 MQuAKE atomic-question recall. SURE is the only row that is best or second-best on erasure and best on locality at once on MCF and MQuAKE, is the best eraser on zsRE (Gen 1.24 against 2.98 for RMU, which halves Spe) with locality kept exactly, and keeps PPL at base on every benchmark, which Remark~\ref{rem:locality} predicts for a method that leaves most retain requests unclassified.

\paragraph{Routing versus residual editing.}
Genie isolates the edit vectors from the learned classifier. On MCF, zsRE and MQuAKE the request-prompt erasure of SURE and Genie is identical because the classifier already fires on every request prompt; the held-out gap therefore measures classifier misses on paraphrases. RWKU differs because its held-out probes ask other facts about the forgotten person rather than the specific trained association, so many such probes have no dedicated edit vector.

% Table 2: read/write layer ablation, Llama-3.2-3B-Instruct, mean over 5 seeds.
\begin{table*}[t]
\centering
\small
\setlength{\tabcolsep}{3.4pt}
\begin{tabular}{l c c ccc ccc ccc ccc}
\toprule
& & & \multicolumn{3}{c}{MCF} & \multicolumn{3}{c}{ZsRE} & \multicolumn{3}{c}{MQuAKE} & \multicolumn{3}{c}{RWKU} \\
\cmidrule(lr){4-6}\cmidrule(lr){7-9}\cmidrule(lr){10-12}\cmidrule(lr){13-15}
Variant & read & write & Eff$\downarrow$ & Gen$\downarrow$ & Spe$\uparrow$ & Eff$\downarrow$ & Gen$\downarrow$ & Spe$\uparrow$ & Eff$\downarrow$ & Gen$\downarrow$ & Spe$\uparrow$ & Eff$\downarrow$ & Gen$\downarrow$ & Spe$\uparrow$ \\
\midrule
Base & -- & -- & 13.46 & 10.26 & 19.88 & 32.37 & 31.20 & 29.79 & 73.23 & 46.44 & 66.13 & 57.60 & 54.77 & 63.70 \\
Ours & 19 & 16 & 0.001 & 0.027 & 19.84 & 0.17 & 1.24 & 29.79 & 0.10 & 1.49 & 55.07 & 2.00 & 33.61 & 61.89 \\
one pass, read late & 19 & 19 & 0.001 & 0.027 & 19.84 & 0.41 & 1.48 & 29.79 & 0.53 & 1.83 & 55.53 & 2.40 & 34.10 & 61.82 \\
one pass, read early & 16 & 16 & 0.001 & 0.013 & 19.80 & 0.04 & 2.76 & 29.79 & 0.17 & 2.62 & 55.25 & 2.40 & 34.81 & 62.18 \\
\bottomrule
\end{tabular}
\caption{Where to read and where to write (Llama-3.2-3B-Instruct, mean over 5 seeds). Reading at block 19 improves held-out generalisation on zsRE relative to reading at block 16 (Gen 1.24 vs 2.76). With the same block-19 reader, writing at block 16 improves MQuAKE generalisation relative to writing at block 19 (Gen 1.49 vs 1.83). Reading at block 19 and writing at block 16 costs three extra blocks for the last prompt token per request and nothing per generated token.}
\label{tab:ablation_layers}
\end{table*}

\paragraph{Where to read, where to write.}
Table~\ref{tab:ablation_layers} moves the read and write blocks while keeping the rest of the method fixed. Reading at block 16 instead of 19 slightly improves the direct zsRE prompt (Eff 0.04 vs 0.17) but hurts held-out paraphrases (Gen 2.76 vs 1.24), suggesting that the later representation is more useful for classifying semantically varied requests. With the same block-19 reader, writing at block 16 improves held-out MQuAKE generalisation relative to writing at block 19 (Gen 1.49 vs 1.83). Thus the read layer and write layer serve different roles: read late for detection, write earlier for a more general edit. The price is three extra blocks for one token per request.

\section{Limitations}
\label{sec:limits}
SURE is behavioural unlearning with runtime components: removing the classifier restores the base model, so it does not satisfy a weights-only threat model, and it does not remove the fact from the model's internal representations; an adversary with access to hidden states is outside the threat model it addresses. The subject gate bounds its scope: a request that never names the subject (a description, a pronoun across turns, a multi-hop chain), or names it in a form the gate does not know (an alias, or a translated place name such as \emph{Wien} for Vienna), is never classified and is answered by the base model. Person names survive translation and the classifier itself is language-agnostic, so cross-lingual questions about people are classified; alias lists per subject would close the remaining gap. Edit vectors are per fact, so entity-level forgetting needs an entity-level classifier. The method is evaluated at 50 facts per split; the cost of rewrites and edit-vector training grows linearly in the number of facts.

\section{Conclusion}
Separating classification from residual editing turns fact-level unlearning into two measurable parts. A late-block classifier decides whether to intervene, an earlier residual vector applies the edit, and every request the classifier declines is answered by the base model. On four benchmarks the combination erases a fact on its request and on held-out phrasings while preserving base-model locality and perplexity; the oracle-classifier comparison and the layer ablation identify whether remaining failures come from the classifier or from the edit itself.

\bibliography{references}
\bibliographystyle{icml2026}

\appendix
\onecolumn
\section{Other models}
\label{app:models}
Tables~\ref{tab:llama8b} and~\ref{tab:qwen4b} repeat Table~\ref{tab:main} on Llama-3.1-8B-Instruct and Qwen3-4B with every baseline. The classifier recipe and edit-vector training are unchanged; only $R$ and $W$ change. $R$ is the block that a sweep of classifier quality over all blocks selects on MCF seed 1 (best F1, earliest block within 0.005 of it with no higher false-positive rate): 20 on Llama-3.1-8B and 26 on Qwen3-4B. On Llama-3.1-8B, $W=18$ is the 3B model's relative write depth; on Qwen3-4B, $W=23$ has the lowest Gen averaged over MCF and zsRE among $W\in\{16,18,21,23,26\}$ in a seed-1 smoke test, with Spe at the same value for all five.

On Llama-3.1-8B the picture of the 3B model repeats: SURE removes nearly all request-prompt recall and most held-out recall with Spe and PPL at the base model's level, it is the only method that is best on locality on MCF, MQuAKE and RWKU while erasing, and Genie on the same edit vectors removes nearly all held-out recall, so the remaining leak is classification. The one-pass variant at block 22 confirms the write-early effect on the request prompt (zsRE Eff 0.13 against 0.08, MQuAKE 1.32 against 0.79), and on RWKU the entity-level Genie retains 50.3 of held-out recall at block 22 against 27.7 at block 18. On Qwen3-4B, SURE keeps perplexity at the base model's level and locality within 0.6 points of it on every benchmark (Spe$^\dagger$ on MQuAKE), and is the best method on zsRE erasure, but its edit vectors are weaker: Genie leaves 9.3 of 30.2 points of zsRE request-prompt recall and 16.1 of 58.8 on MQuAKE, so on this model the remaining leak on these two benchmarks is editing, not classification.

\begin{table*}[t]
\centering
\small
\setlength{\tabcolsep}{3.2pt}
\begin{tabular}{l cccc cccc cccc cccc}
\toprule
& \multicolumn{4}{c}{MCF} & \multicolumn{4}{c}{ZsRE} & \multicolumn{4}{c}{MQuAKE} & \multicolumn{4}{c}{RWKU} \\
\cmidrule(lr){2-5}\cmidrule(lr){6-9}\cmidrule(lr){10-13}\cmidrule(lr){14-17}
Method & Eff$\downarrow$ & Gen$\downarrow$ & Spe$\uparrow$ & PPL & Eff$\downarrow$ & Gen$\downarrow$ & Spe$\uparrow$ & PPL & Eff$\downarrow$ & Gen$\downarrow$ & Spe$\uparrow$ & PPL & Eff$\downarrow$ & Gen$\downarrow$ & Spe$\uparrow$ & PPL \\
\midrule
Base        & 17.09 & 13.01 & 22.48 & 11.03 & 39.81 & 36.44 & 30.56 & 11.03 & 72.82 & 63.48 & 67.76 & 11.03 & 78.40 & 83.98 & 84.71 & 11.03 \\
\midrule
GA          & 0.019 & 0.017 & 0.00 & 88.85 & 23.86 & 23.47 & 25.97 & 13.26 & 2.38 & \underline{2.28} & 3.82 & 2.0e5 & 15.20 & 16.71 & 14.83 & 12.42 \\
GD          & 0.007 & \textbf{0.008} & 0.00 & 74.78 & 24.01 & 23.67 & 25.97 & 13.89 & 3.12 & 2.45 & 4.62 & 1969 & 15.60 & 17.42 & 15.68 & 12.72 \\
NPO         & 0.012 & 0.014 & 0.00 & 95.48 & 21.43 & 21.20 & 25.02 & 13.29 & 2.35 & 2.43 & 3.67 & 175.10 & 5.60 & 11.56 & 10.40 & 12.64 \\
RMU         & \underline{0.005} & \underline{0.013} & 0.60 & 10.73 & \underline{1.13} & \textbf{0.93} & 6.85 & 10.74 & \textbf{0.00} & \textbf{0.00} & 0.10 & 10.78 & \textbf{0.00} & \textbf{0.00} & 0.00 & 10.85 \\
ROME        & 3.595 & 4.362 & 20.00 & 11.12 & 34.64 & 34.88 & 30.53 & 11.10 & 7.88 & 17.18 & 32.05 & 11.19 & 6.40 & \underline{7.27} & 51.81 & 11.08 \\
MEMIT       & 6.316 & 7.731 & \underline{21.48} & \underline{11.09} & 37.65 & 37.27 & \textbf{30.82} & \underline{11.07} & 28.02 & 27.84 & 55.17 & \underline{11.16} & 53.20 & 62.22 & \underline{80.24} & \underline{11.07} \\
ZeroUnlearn & 1.932 & 4.477 & 18.32 & 11.17 & 33.77 & 32.52 & 30.43 & 11.18 & 24.96 & 28.26 & \underline{55.30} & 11.20 & 33.20 & 46.48 & 75.04 & 11.15 \\
\midrule
\textbf{Ours}& \textbf{0.000} & 0.353 & \textbf{22.44} & \textbf{11.03} & \textbf{0.08} & \underline{3.27} & \underline{30.56} & \textbf{11.03} & \underline{0.79} & 4.95 & \textbf{58.41} & \textbf{11.03} & \underline{3.60} & 49.21 & \textbf{82.62} & \textbf{11.03} \\
\quad one pass, block 22 & 0.000 & 0.353 & 22.44 & 11.03 & 0.13 & 4.09 & 30.57 & 11.03 & 1.32 & 4.94 & 58.49 & 11.03 & 3.60 & 43.94 & 82.44 & 11.03 \\
Genie       & 0.000 & 0.002 & 22.48 & 11.03 & 0.08 & 0.08 & 30.56 & 11.03 & 0.79 & 0.79 & 58.75 & 11.03 & 1.60 & 83.98 & 84.71 & 11.03 \\
Genie (entity)& -- & -- & -- & -- & -- & -- & -- & -- & -- & -- & -- & -- & 1.60 & 27.65 & 84.71 & 11.03 \\
\bottomrule
\end{tabular}
\caption{Llama-3.1-8B-Instruct (32 blocks), mean over 5 seeds. Classifier at block 20, edit vectors at block 18; ``one pass'': both at block 22. GA, GD and NPO train the MLP down projections; RMU steers block 8. On MQuAKE, Spe$^\dagger$ is 70.10 (Base) and 69.69 (Ours). Bold and underline as in Table~\ref{tab:main}. Standard deviations in \texttt{RESULTS\_PAPER\_TABLES.md}.}
\label{tab:llama8b}
\end{table*}

\begin{table*}[t]
\centering
\small
\setlength{\tabcolsep}{3.2pt}
\begin{tabular}{l cccc cccc cccc cccc}
\toprule
& \multicolumn{4}{c}{MCF} & \multicolumn{4}{c}{ZsRE} & \multicolumn{4}{c}{MQuAKE} & \multicolumn{4}{c}{RWKU} \\
\cmidrule(lr){2-5}\cmidrule(lr){6-9}\cmidrule(lr){10-13}\cmidrule(lr){14-17}
Method & Eff$\downarrow$ & Gen$\downarrow$ & Spe$\uparrow$ & PPL & Eff$\downarrow$ & Gen$\downarrow$ & Spe$\uparrow$ & PPL & Eff$\downarrow$ & Gen$\downarrow$ & Spe$\uparrow$ & PPL & Eff$\downarrow$ & Gen$\downarrow$ & Spe$\uparrow$ & PPL \\
\midrule
Base        & 11.13 & 9.78 & 12.64 & 14.07 & 30.21 & 30.37 & 35.59 & 14.07 & 58.82 & 24.91 & 51.85 & 14.07 & 26.00 & 30.43 & 24.89 & 14.07 \\
\midrule
GA          & 0.018 & 0.788 & 0.72 & 650.50 & 21.85 & 21.18 & 33.54 & 17.13 & 11.00 & 15.66 & 16.81 & 80.77 & 12.40 & 15.15 & 12.59 & 14.74 \\
GD          & 0.032 & 0.786 & 0.80 & 470.91 & 21.98 & 21.21 & 33.32 & 18.29 & 11.19 & 15.94 & 17.16 & 99.91 & 12.40 & 15.28 & 12.81 & 15.05 \\
NPO         & 0.128 & 0.382 & 0.76 & 175.93 & 21.38 & 20.27 & 33.90 & 17.53 & 11.31 & 16.47 & 17.08 & 60.55 & 13.20 & 15.74 & 12.38 & 15.99 \\
RMU         & 0.027 & \underline{0.044} & 1.00 & 13.67 & \underline{10.33} & \underline{11.88} & 27.37 & 13.99 & \underline{1.81} & \underline{2.25} & 5.35 & \underline{14.04} & \textbf{0.80} & \textbf{0.36} & 0.07 & 13.94 \\
ROME        & \textbf{0.000} & \textbf{0.003} & 0.00 & 2140 & 14.82 & 14.90 & 25.06 & 15.93 & \textbf{0.00} & \textbf{0.00} & 0.01 & 8.2e6 & \underline{2.80} & \underline{4.09} & 14.62 & 14.10 \\
MEMIT       & 0.743 & 2.333 & \underline{12.44} & \underline{14.21} & 27.45 & 26.98 & \underline{35.04} & \underline{14.14} & 26.95 & 20.96 & \textbf{48.16} & 14.25 & 18.40 & 15.22 & \underline{23.07} & \underline{14.09} \\
ZeroUnlearn & 0.043 & 0.590 & 9.72 & 15.24 & 24.73 & 25.39 & 34.54 & 15.36 & 20.43 & 18.85 & 37.51 & 17.52 & 9.60 & 13.55 & 21.04 & 14.46 \\
\midrule
\textbf{Ours}& \underline{0.003} & 0.118 & \textbf{12.52} & \textbf{14.07} & \textbf{9.33} & \textbf{11.22} & \textbf{35.59} & \textbf{14.07} & 16.05 & 15.15 & \underline{45.98} & \textbf{14.07} & 4.00 & 20.71 & \textbf{24.53} & \textbf{14.07} \\
Genie       & 0.003 & 0.117 & 12.64 & 14.07 & 9.33 & 9.43 & 35.59 & 14.07 & 16.05 & 14.35 & 46.44 & 14.07 & 2.00 & 30.43 & 24.89 & 14.07 \\
Genie (entity)& -- & -- & -- & -- & -- & -- & -- & -- & -- & -- & -- & -- & 2.00 & 20.94 & 24.89 & 14.07 \\
\bottomrule
\end{tabular}
\caption{Qwen3-4B (36 blocks), mean over 5 seeds. Classifier at block 26, edit vectors at block 23. RMU steers block 9. ROME diverges on MCF and MQuAKE on this model (PPL $2\cdot10^3$ and $8\cdot10^6$). On MQuAKE, Spe$^\dagger$ is 52.13 (Base) and 51.58 (Ours). Bold and underline as in Table~\ref{tab:main}.}
\label{tab:qwen4b}
\end{table*}

\section{RWKU with an entity-level classifier}
\label{app:rwku}
The RWKU row of Table~\ref{tab:main} uses the fact-level classifier. With the entity-level gate (fire on any probe that names a protected person; the heads choose which of the person's edit vectors), held-out recovery can be compared with the corresponding oracle bound. \todo{insert outputs/l16\_ours\_rwku\_subject\_v1}

\section{Classifier training data}
\label{app:classifier}
Per fact: 20 positives (the direct prompt; base-model rewrites that name the subject, omit the answer, are novel, and are answered by the base model with the fact's answer; up to 8 context-prefix copies), 20 negatives (same-subject competitor facts up to half, model-written hard negatives, subject transplants; dropped on answer collision). A third of the rewrites is held out for bias calibration at 98\% recall. The classifier input $\mathbf u(x)$ (Equation~\eqref{eq:feature}) is the unit-normalised, mean-centred block-$R$ state at the last prompt token; no dimensionality reduction.

\section{Edit-vector training}
\label{app:edits}
Five rounds; each round regenerates greedy continuations (16 tokens) of six training views per fact under the current edit vectors and adds them to a replay buffer, then 150 Adam steps (learning rate $0.01$, batch 8) on Equation~\eqref{eq:loss} with $\tau=10^{-4}$, $\lambda=0.1$, each step followed by projection onto $\|\boldsymbol\delta_i\|_2\le\rho_W$ (the median block-$W$ residual norm); an association the classifier never selects during training is trained with its edit vector forced on. Teacher-forced gold prefixes are added as views on zsRE, MQuAKE and RWKU, whose metrics score every answer token under the gold prefix.

\end{document}
